\documentclass{article}
\usepackage{graphicx, amsmath} % Required for inserting images

\title{Apostol, Chapter 1 - Exercise 16}
\author{Afonso Vilhena Francisco}
\date{31 August 2026}

\begin{document}

\maketitle

$\mathbf{16.}$ \textbf{Prove that if} $\mathbf{2^n-1}$ \textbf{is prime, then} $\mathbf{n}$ \textbf{is prime.}

\vspace{0.5cm}

Suppose the $n$ is composite, that is, $n=ab$, for some integers $a,b>1.$

For any $j>1$, we have that $$x^j-1=(x-1)(x^{j-1}+x^{j-2}+...+x+1).$$

Taking $x=2^a$ and $j=b$ we get

\begin{align*}
    2^n-1 = 2^{ab}-1 = (2^a-1)(2^{a(b-1)}+2^{a(b-2)}+...+2^a+1)
\end{align*}

Since $a>1$, we have that $2^a-1>1$ and $2^{a(b-1)}+2^{a(b-2)}+...+2^a+1$. Thus, $2^n-1$ is not a prime (Contradiction).

Therefore, $n$ is prime.

\end{document}
